% ═══════════════════════════════════════════════════════════════════════════
\chapter{Introducción}
% ═══════════════════════════════════════════════════════════════════════════

El market making constituye uno de los problemas centrales de la microestructura de
mercados modernos. Un agente que actúa como market maker provee liquidez continua al
mercado publicando simultáneamente cotizaciones de compra y venta, buscando capturar
la diferencia entre ambas ---el spread--- a cambio de asumir el riesgo de
inventario y el riesgo de selección adversa \citep{GlostenMilgrom1985}.

Con el auge de los mercados electrónicos y el trading algorítmico, el diseño de
estrategias óptimas de market making ha adquirido una relevancia tanto académica como
práctica considerable. El modelo seminal de \citet{AvellanedaStoikov2008} estableció
el marco de referencia: formulando el problema como uno de control estocástico en
tiempo continuo, los autores derivan cotizaciones óptimas en función del inventario
acumulado y la volatilidad del activo. Sin embargo, este modelo descansa sobre una
especificación restrictiva: las ejecuciones se describen mediante procesos de Poisson
condicionales con forma exponencial en la distancia y parámetros fijos, sin depender
explícitamente de covariables contemporáneas del LOB ni de la historia reciente de
eventos.

La evidencia empírica muestra las limitaciones de esta especificación. Los datos de
alta frecuencia muestran que la dinámica de llegada de órdenes exhibe
\textit{clustering} temporal, dependencias entre eventos consecutivos y una fuerte
sensibilidad al estado observable del libro de órdenes, fenómenos que los procesos
de Poisson con parámetros fijos no pueden reproducir \citep{BacryHawkes2015}. Esta
brecha entre el modelo y la realidad puede producir intensidades mal calibradas
y afectar tanto la política de cotización como el control del riesgo de inventario.

Este trabajo propone cerrar esa brecha integrando herramientas de aprendizaje
estadístico dentro del marco de control estocástico, mediante un modelo
semiparamétrico que estima las intensidades de órdenes como funciones del estado
observable del mercado. El caso de estudio empírico se desarrolla sobre BTC/USDT
en Binance Spot, un mercado de criptoactivos que opera de manera continua y sobre el
cual se implementó una captura propia de datos históricos de alta frecuencia
directamente desde la infraestructura pública de mercado de Binance.


El trabajo responde dos preguntas que se mantienen separadas a lo largo del documento. La primera es predictiva: cuánto mejora la propuesta la estimación de la intensidad de llegada de órdenes agresivas a cada distancia del precio medio, frente a Avellaneda--Stoikov y a un Hawkes bivariado que reciben las mismas mejoras aplicables. La segunda es económica: qué valor produce esa mejora cuando se traduce en una política de cotización, bajo una regla de ejecución, unos costos y un inventario comunes. Una contribución predictiva puede coexistir con un resultado económico negativo, siempre que ambos se sustenten con el mismo rigor. Eso es lo que se encuentra aquí.

Las contribuciones son cuatro:
\begin{enumerate}[nosep]
\item una especificación semiparamétrica de la intensidad que incorpora el estado del libro, la periodicidad, una curva monótona de distancia modulada por el estado y covariables de barridos;
\item una comparación con referencias reforzadas, en la que E1 y E2 reciben los mismos hallazgos que E3, evaluada con \textit{walk-forward} y con especificaciones finales congeladas sobre un periodo de test;
\item la resolución numérica de la ecuación HJB con intensidades dependientes del estado y su evaluación económica con la misma regla de ejecución para todos los modelos;
\item una explicación cuantitativa de por qué la mejora predictiva no se traduce en rentabilidad en BTC/USDT, apoyada en la literatura reciente y en una comparación con tres motores de simulación limitada a configuraciones anteriores del entrenamiento.
\end{enumerate}

El Capítulo~\ref{chap:marco-teorico} presenta el marco teórico, el Capítulo~\ref{chap:metodologia} la metodología, el Capítulo~\ref{chap:datos} la captura y reconstrucción de los datos y el Capítulo~\ref{chap:resultados} los resultados, primero los predictivos y después los económicos. El Capítulo~\ref{chap:discusion} los interpreta frente a la literatura y declara las limitaciones, y el Capítulo~\ref{chap:conclusiones} concluye.
